\subsection{From VAEs to VQ-VAEs}
\begin{frame}{}
	\LARGE VAE: \textbf{Two Problems on the Road to VQ-VAE}
\end{frame}

\begin{frame}{Problem 1: The Averaging Problem}
\begin{itemize}
    \item The decoder must produce \textbf{one} output for each $z$, and the Gaussian likelihood (MSE) makes the optimal output the \textbf{average} of everything that $z$ could correspond to.
    \item The KL term makes this worse: it deliberately \textbf{overlaps} the posteriors of different inputs, so one latent region is shared by several training images and decodes to their compromise.
    \item Result: the characteristic \textbf{blurriness} of VAE samples; sharp textures and fine details are averaged away.
    \item Perceptual losses treat the symptom; the multimodality itself calls for a more expressive latent structure.
\end{itemize}
\end{frame}

\begin{frame}{Problem 2: The Prior Hole Problem}
\begin{itemize}
    \item Training only ever feeds the decoder latents drawn from the posteriors $q(z \mid x)$ of \textbf{real data}; generation samples from the \textbf{prior} $p(z) = \mathcal{N}(0, I)$.
    \item The KL term pushes each posterior toward the prior, but the \textbf{aggregate} of all posteriors never matches it exactly: some prior regions receive \emph{no} training data. These are \textbf{prior holes}.
    \item A sample landing in a hole is decoded by a network that has never seen that region: outputs degrade or break entirely (the failure our latent-space figure showed for plain autoencoders, in a milder form).
    \item Related failure: \textbf{posterior collapse}; with a powerful decoder, the model can ignore $z$ altogether and the latent space carries no information.
    \item Ways out: richer learned priors, hierarchical VAEs, or a different idea entirely: \textbf{discrete} latent codes with a \textbf{learned prior}. That is VQ-VAE.
\end{itemize}
\end{frame}
